% Generated by `python -m analysis.export_figures`. Copy figure blocks into report.tex.
% figure* spans both columns; change to figure and width=\linewidth for one column.

\begin{figure*}[t]
  \centering
  \includegraphics[width=\textwidth]{county_coverage_map}
  \caption{86\% of the 271 mapped sites are east of 106.75\textdegree{} W, along the Rio Grande side; only 17 are within 25 km of Cuba, and none of those is a monitoring well in the collected data. Every mapped site in the project's datasets (1926-2026 records), colored by what it measures; outline = Sandoval County. A dot is a site with any record in the collected files, not necessarily active monitoring: many are one-time tests from decades ago (hover for record years). Water-level records = USGS monitoring wells and NMBGMR wells; well construction and aquifer tests = OSE aquifer tests and well logs; surface-water stations = USGS stream, canal, and reservoir sites, including Rio Puerco at Cuba. The thin vertical line at 106.75\textdegree{} W is the east/west split used in the title. 3 sites lie just outside the simplified county line but are in the downloaded county files. One OSE well log isn't mapped because its coordinates disagree with its legal description. Town locations are approximate; in the app, zoom in to separate the dense Rio Grande cluster. Scope: Sandoval County, including the Cuba area. Source: data/clean/groundwater\_levels\_monthly.csv, data/clean/aquifer\_tests.csv, data/clean/groundwater\_levels\_nmbgmr.csv, data/clean/well\_logs.csv, data/water\_flow\_data\_sandoval\_county\_sources.csv (dataset 1), data/clean/streamflow\_measurements.csv; county outline: data/sandoval\_county\_boundary.geojson (US Census Bureau, via plotly/datasets).}
  \label{fig:county_coverage_map}
\end{figure*}

\begin{figure*}[t]
  \centering
  \includegraphics[width=\textwidth]{groundwater_hydrographs}
  \caption{5 of 9 monitored wells fell over their records and 3 rose; the Sandia site has risen 13 ft since 2007. Monthly mean depth to water at the 9 USGS monitoring wells (4 sites), 1998-2026. Axis flipped: up = higher water table. Each panel is one monitoring site. Wells marked A or B are at the same spot but measured at a different depth, so they are drawn separately; at Lincoln Middle School and the Sandia site, Well and Well A nearly coincide (Well A is dashed). Every panel spans the same 37 ft, so steepness can be compared across panels. Lines break where months are missing. These are water-level indicators at 4 sites, not measures of how much water is left. Scope: Sandoval County, not Cuba: 4 sites between Rio Rancho and Cochiti, 65-95 km from Cuba. Source: USGS Water Data statistics, parameter 72019, cleaned by J. Wauneka (data/clean/groundwater\_levels\_monthly.csv).}
  \label{fig:groundwater_hydrographs}
\end{figure*}

\begin{figure*}[t]
  \centering
  \includegraphics[width=\textwidth]{groundwater_trends}
  \caption{In 2016-2025, 4 of 7 wells were still falling (Lincoln Middle School and Dome Road well); only the Sandia site wells rose. Sen's slope of annual mean water level per well: whole record (circles) vs 2016-2025 (diamonds); whiskers = 95\% interval; gray years = each well's record. Trend = Sen's slope of annual means (years with at least 9 months of data). Sen's slope assumes one direction, so a record that fell and then rose, like the Sandia site (see the hydrographs), gets one average slope. Color shows the Mann-Kendall test: red = falling, blue = rising, gray = no clear trend (p $\geq$ 0.05). Sign is flipped from depth so that negative = water table falling. Both the p-values and the intervals assume each year is independent; groundwater years aren't (a low year tends to follow a low year), so both look more certain than they are. Records differ in length (see the years by each well). Scope: Sandoval County, not Cuba: the same 9 USGS wells, 65-95 km from Cuba. Source: USGS Water Data statistics, parameter 72019 (data/clean/groundwater\_levels\_monthly.csv); annual means, Sen's slope, Mann-Kendall test (analysis/stats.py).}
  \label{fig:groundwater_trends}
\end{figure*}

\begin{figure*}[t]
  \centering
  \includegraphics[width=\textwidth]{placitas_remeasured}
  \caption{Of 5 Placitas wells measured in 1997-1998 and 2018-2019, 3 changed $\leq$5 ft; PC-014 is 51 ft deeper and PC-025 57 ft shallower (unstable baseline). Average depth to water per period for NMBGMR wells measured in both periods; shaded bar = range of the 1997-1998 readings; n = readings per period. The 1997-1998 readings were taken about every 63 days (median gap), so the shaded bar shows how much a well varied then; a change smaller than that range may just be seasonal. \dag{} = at least one 1997-1998 reading carries LevelStatus R (in USGS coding, 'recently pumped'), which can make the water level read deeper than normal. The measuring method also changed between periods (code T in 1997-1998, S/T in 2018-2019; T is electric tape and S steel tape in USGS coding). Wells measured in only one period (82 of 87) can't show change and are left out. Scope: Placitas area of Sandoval County, about 90 km from Cuba (not Cuba). Source: NM Bureau of Geology and Mineral Resources groundwater levels (data/clean/groundwater\_levels\_nmbgmr.csv, prepared in analysis/data.py).}
  \label{fig:placitas_remeasured}
\end{figure*}

\begin{figure*}[t]
  \centering
  \includegraphics[width=\textwidth]{well_depth_vs_water}
  \caption{Where water is deeper, wells are drilled deeper ($\rho$ = 0.63-0.64 within the 2 larger record types), but half of 180 wells have water 41-241 ft down. One point per well: well depth vs depth to water, from three record types, 1930-2022. Points far above the diagonal have a lot of water standing above the well bottom; points near it have little. Colors follow the project's data types (blue = water-level records, orange = well construction and tests). The three sources measure water differently (a 1997-2019 measurement, the level on the day of an aquifer test, the depth where water was first hit while drilling), so treat them as context, not one comparable series. 1 record was left out because its water level is deeper than the well itself (RG-54556x: 30 ft well, water at 458 ft), likely a data-entry error. Scope: Sandoval County, including the Cuba area (most wells are around Placitas and Cochiti; 7 within 25 km of Cuba). Source: NMBGMR groundwater levels (data/clean/groundwater\_levels\_nmbgmr.csv), NM OSE aquifer tests (data/clean/aquifer\_tests.csv), NM OSE well logs (data/clean/well\_logs.csv).}
  \label{fig:well_depth_vs_water}
\end{figure*}

\begin{figure*}[t]
  \centering
  \includegraphics[width=\textwidth]{depth_by_area}
  \caption{Among recorded wells, depth to water was smallest around Cochiti-Pe\~{n}a Blanca (median 22 ft) and largest around Bernalillo-Rio Rancho (median 375 ft; 9 of 11 are aquifer tests from 1987-2007). Depth to water for each well with a recorded level, 1930-2022, grouped by the nearest town; black tick = area median. Areas are a simple screening grouping: each well goes to the nearest of seven towns within 20 km (25 km for Cuba, to match the project's Cuba area), otherwise 'Elsewhere in county'. They are not hydrogeologic regions. One point per well, including wells without a recorded well depth; medians are shown only where an area has at least 5 wells. 1 well without a usable location is left out. Scope: Sandoval County, including the Cuba area. Source: NMBGMR groundwater levels (data/clean/groundwater\_levels\_nmbgmr.csv), NM OSE aquifer tests (data/clean/aquifer\_tests.csv), NM OSE well logs (data/clean/well\_logs.csv).}
  \label{fig:depth_by_area}
\end{figure*}

\begin{figure*}[t]
  \centering
  \includegraphics[width=\textwidth]{formation_properties__transmissivity_drawdown_ft2_per_day}
  \caption{Transmissivity varies widely inside formations: in 3 of 9, values span a wider range than the formation medians do (450-fold), up to 61,000-fold in the Madera Limestone. Transmissivity, drawdown method (ft\textsuperscript{2}/day) from each aquifer test, by geologic formation (formations with at least 3 values, plus the two Cuba-area formations); black tick = median of wells; dotted line = median of all wells. Transmissivity is how easily water moves through the full thickness of the aquifer; higher means a more productive aquifer. Each dot is one reported value; a well tested more than once (or reported on two rows for the same day) has more than one dot, so the gray label counts values and wells. Medians are taken per well first, so a well tested twice counts once, and are drawn only for formations with at least 3 tested wells. Wells without a recorded formation (22 of the 108 wells with a value) count only in the all-wells median. Tests span 1930-2008. Scope: Sandoval County, including the Cuba area. Source: NM Office of the State Engineer aquifer tests, cleaned by J. Wauneka (data/clean/aquifer\_tests.csv).}
  \label{fig:formation_properties}
\end{figure*}

\begin{figure*}[t]
  \centering
  \includegraphics[width=\textwidth]{cuba_area_map}
  \caption{Within 25 km of Cuba: 15 wells with one-time aquifer tests (1957-1979), 1 well log, and 1 stream gage; in the collected data, no well there is measured repeatedly. Sites within 60 km of Cuba, colored by what they measure; rings at 10, 25, and 50 km. Most sites near Cuba are aquifer tests done once, decades ago; they give well depth, aquifer properties, and the water level on test day, not a trend. The Rio Puerco gage measures the river, not groundwater. The 25 km 'Cuba area' radius is an analysis choice, not an official boundary. Town locations are approximate. Scope: Cuba and surroundings (within 60 km). Source: data/clean/groundwater\_levels\_monthly.csv, data/clean/aquifer\_tests.csv, data/clean/groundwater\_levels\_nmbgmr.csv, data/clean/well\_logs.csv, data/water\_flow\_data\_sandoval\_county\_sources.csv (dataset 1), data/clean/streamflow\_measurements.csv; county outline: data/sandoval\_county\_boundary.geojson (US Census Bureau, via plotly/datasets).}
  \label{fig:cuba_area_map}
\end{figure*}

\begin{figure*}[t]
  \centering
  \includegraphics[width=\textwidth]{cuba_nearest_records}
  \caption{The nearest well with repeated water-level measurements is 38 km from Cuba; within 25 km there are only one-time tests at 15 wells (1957-1979), 1 well log, and 1 stream gage. For each dataset, the distance from Cuba's village center to its closest site, and that site's record years. Distances are straight-line from the approximate village center. A short bar only means a site is close, not that it answers the question: the closest records are one-time aquifer tests and river visits, neither of which tracks groundwater over time. Scope: Cuba and the nearest site in each dataset. Source: data/clean/groundwater\_levels\_monthly.csv, data/clean/aquifer\_tests.csv, data/clean/groundwater\_levels\_nmbgmr.csv, data/clean/well\_logs.csv, data/water\_flow\_data\_sandoval\_county\_sources.csv (dataset 1), data/clean/streamflow\_measurements.csv.}
  \label{fig:cuba_nearest_records}
\end{figure*}

\begin{figure*}[t]
  \centering
  \includegraphics[width=\textwidth]{ojo_alamo_depth_map}
  \caption{On the 1990 USGS map, Cuba sits at the southeastern end of the Ojo Alamo outcrop, where the sandstone is shallow; its top deepens northward past 1,000 and 3,000 ft. Crop of USGS Hydrologic Atlas HA-720-B, Figure 5 (approximate depth to the top of the Ojo Alamo Sandstone), with labels added. The map itself is unchanged; only the white labels and arrows are added. Brown is where the Ojo Alamo Sandstone crops out at the surface; each red line joins points where its top is the same depth below land surface, every 500 ft. 'Cuba' is the map's own town label. Dataset 5 is a scanned plate, so values are read by eye and nothing here is digitized. Scope: Cuba area and the eastern San Juan Basin (regional map, 1990). Source: USGS Hydrologic Investigations Atlas HA-720-B (Thorn et al. 1990), Figure 5 (data/data5\_ojo\_alamo\_sandstone\_pub.pdf); public domain.}
  \label{fig:ojo_alamo_depth_map}
\end{figure*}

\begin{figure*}[t]
  \centering
  \includegraphics[width=\textwidth]{ojo_alamo_vs_san_jose}
  \caption{Only 8 test events at 7-8 wells name these formations (7 Ojo Alamo, 1 San Jose; 1953-1977): too few to describe either formation, let alone compare them. Every aquifer test recorded in the Ojo Alamo Sandstone or San Jose Formation (one dot per reported value), with published San Jose transmissivity for comparison. A test event is one well on one date; some events are reported on two rows (e.g. separate drawdown results), so n counts values. Two Ojo Alamo well names, 'Cuba Independent Rural School Well No 1' and '21.1W.28.143', share a location and depth and may be the same well. Open violet diamonds are published San Jose Formation transmissivity compiled in Blandford et al. (2025, Table 1): well SJ-2259, screened in the Cuba Mesa Member near Lindrith (estimated from specific capacity, assuming 70\% well efficiency), and a range from two tests in Stone et al. (1983), location not given; both converted from m\textsuperscript{2}/day. A value entered twice for the same test is drawn once. Depth to water and the two right panels use log scales. Blank panels mean the value wasn't recorded. The tests date from 1953-1977 and don't show which formation Cuba's municipal wells use today. Scope: Sandoval County, including the Cuba area: tests 1-45 km from Cuba; published values from the San Juan Basin. Source: NM OSE aquifer tests (data/clean/aquifer\_tests.csv); published values: Blandford et al. (2025), NMGS Guidebook 75, Table 1 (data/data6\_cuba\_mesa\_member\_of\_san\_jose\_formation.pdf).}
  \label{fig:ojo_alamo_vs_san_jose}
\end{figure*}

\begin{figure*}[t]
  \centering
  \includegraphics[width=\textwidth]{rio_puerco_visits}
  \caption{Rio Puerco at Cuba had measurable flow on all 4 visits in Apr-May and on none of the 6 visits from Aug to Mar. Each dot is one USGS field visit to station 08332525 (2019-2026); left by date (labeled with the month), right by month of year; not a continuous record. No line is drawn between visits because the river was not measured in between. 'No measurable flow' is what the visit recorded at that moment (the channel may still hold pools), not that the river was dry all season. This station measures the river at Cuba, not Cuba's groundwater or water supply. Scope: Cuba (measured at Cuba), surface water only (gage 2.3 km from the village center). Source: USGS station 08332525, Rio Puerco at Cuba, NM (data/clean/streamflow\_measurements.csv).}
  \label{fig:rio_puerco_visits}
\end{figure*}

\begin{figure*}[t]
  \centering
  \includegraphics[width=\textwidth]{coverage_timeline}
  \caption{Groundwater levels are monitored every year only since 1999, at 4 sites in the east of the county; the collected data have no water-level monitoring in the Cuba area. Number of sites with at least one record, by year (1926-2026); top: whole county, bottom: the Cuba area only. Darker = more sites recorded that year. The USGS row counts monitoring sites (locations; the 9 wells sit at 4 sites). Aquifer tests count wells tested that year (each test also records a one-time water level, which isn't monitoring). Streamflow includes the Rio Puerco at Cuba visits. Gaps are years with no record in these files; other agencies may have data. Scope: Whole project (all datasets). Source: data/clean/groundwater\_levels\_monthly.csv, data/clean/aquifer\_tests.csv, data/clean/groundwater\_levels\_nmbgmr.csv, data/clean/well\_logs.csv, data/water\_flow\_data\_sandoval\_county\_sources.csv (dataset 1), data/clean/streamflow\_measurements.csv.}
  \label{fig:coverage_timeline}
\end{figure*}

\begin{figure*}[t]
  \centering
  \includegraphics[width=\textwidth]{evidence_matrix}
  \caption{Per the dataset briefs, no dataset gives more than limited support for research question 3, about groundwater near Cuba; coverage, trend, and well questions are the best supported. How well each dataset can answer each research question; blank = not applicable or not assessed. Ratings come from Dylan's one-page dataset briefs; each cell shows a short form of the brief's wording (e.g. 'Best fit' for 'Strongest fit (partial)'), grouped into four colors: None = None; Limited = Limited, Partial, Potential, Conditional, Context; Moderate = Moderate, Best fit, Useful; Strong = Strong. The briefs rated the proposal's nine questions; their ratings are shown under the team's six (questions 1-4 unchanged, 8 $\rightarrow$ 5, 9 $\rightarrow$ 6; the three dropped questions are left out). \dag{} = no brief; rated by the analysis (the river part of USGS dataset 1 and the OSE aquifer tests). \ddag{} = an image-only map plate: it is shown as an attributed map on the Cuba page but not digitized, so its numbers aren't charted, and its 'Strong' for question 4 holds only once it is digitized. This is a judgment summary of what each dataset could support, not a measurement. Scope: Whole project (all datasets). Source: Dataset briefs 1-7 (docs/files/support\_doc1-7) by D. Gonzales; \dag{} no brief, rated by the analysis; \ddag{} image-only map plate, not digitized.}
  \label{fig:evidence_matrix}
\end{figure*}

\begin{figure*}[t]
  \centering
  \includegraphics[width=\textwidth]{aquifer_sensitivity_classes}
  \caption{High sensitivity covers 25\% of the 613 mapped areas but 48\% of the listed area; without locations, none of it can be placed on the map or linked to Cuba. Aquifer-sensitivity classes in the downloaded table (dataset 3): share of polygons and share of listed area. Area shares use the table's area\_sq\_deg column, in square degrees (the source doesn't state units; square degrees shrink slightly northward), so the shares are approximate. Polygon counts and area shares differ because a few large polygons dominate. This describes the downloaded inventory, not Cuba. Scope: Probably statewide, not Sandoval County or Cuba (the file has no locations; if its area units are square degrees, the listed areas add up to about the size of New Mexico). Source: New Mexico Water Data, aquifer sensitivity (data/clean/aquifer\_sensitivity.csv).}
  \label{fig:aquifer_sensitivity_classes}
\end{figure*}
